\section{The single-link law of SU(2)}\label{app:sharp}

This appendix proves the sharp quarter, $W_1(\nu_b,\nu_{b'})\le\frac14|b-b'|$
for the one-link law of $\SU(2)$ in the chordal metric, which
the proof of Theorem~\ref{thm:A} for $\SU(2)$, Lemma~\ref{lem:onelink}(i) and the whole of
Sec.~\ref{sec:stars} use, together with the exact value $\frac14$ at $b=0$,
which shows that the single-link argument of Sec.~\ref{sec:state} stops at
$\beta_W=\frac13$. The argument is that of~\cite{BB26a}; it is given in full
so that the paper is self-contained. The computations it calls on are in
Appendix~\ref{app:vmfcomp}, and the facts about Bessel functions in
Appendix~\ref{app:bessel}.

\subsection{The von Mises--Fisher law}

Identify $\SU(2)$ with the unit quaternions, that is with the unit sphere
$S^3\subset\R^4$, so that $\ReTr U=2x_0(U)$ and $\ReTr(UV)=2\ip u{\bar v}$
for the Euclidean inner product of $\R^4$ and quaternion conjugation
$v\mapsto\bar v$. Fix a link $e$ and the field off $e$. Every plaquette
through $e$ can be rotated cyclically so that its trace reads $\ReTr(Q_eA_p)$
with $A_p\in\SU(2)$, the staple of $e$ in $p$. The action therefore depends on
$Q_e$ only through
\[
N\beta\sum_{p\ni e}\ReTr(Q_eA_p)=\ip{x}{b},\qquad
b=4\beta\sum_{p\ni e}\bar A_p\in\R^4,
\]
with $x=Q_e\in S^3$. Since there are six staples, each of norm one,
\begin{equation}\label{q-eq:kappamax}
\kappa:=|b|\le y:=24|\beta| .
\end{equation}
The conditional law of $Q_e$ is the von Mises--Fisher law\footnote{On the
sphere $S^{p-1}\subset\R^p$ the von Mises--Fisher law with mean direction
$\mu$ and concentration $\kappa$ has density proportional to
$e^{\kappa\ip x\mu}$ with respect to the surface measure, and its normalising
constant is $(2\pi)^{p/2}I_{p/2-1}(\kappa)/\kappa^{p/2-1}$; it is the maximum
entropy law with prescribed mean resultant, and it plays on the sphere the
role of the Gaussian in $\R^p$. For $p=4$ the Bessel index is $1$. The mean
resultant length $\E\ip x\mu=I_{p/2}(\kappa)/I_{p/2-1}(\kappa)$ is the
function $A$ of Lemma~\ref{q-lem:moments}. See Mardia and Jupp~\cite{MJ00},
Ch.~9, for these facts and for the tangent-normal decomposition used in
Appendix~\ref{app:vmfcomp}.}
\begin{equation}\label{q-eq:vmf}
\nu_b(dx)=Z(\kappa)^{-1}e^{\ip xb}\,\sigma(dx)
\end{equation}
on the unit sphere, with $\sigma$ the normalised surface measure. We use
coordinates in which $b=\kappa e_1$. The density of $t=x_1$ is proportional
to $(1-t^2)^{1/2}e^{\kappa t}$ on $[-1,1]$, so $Z(\kappa)=2I_1(\kappa)/\kappa$
with the modified Bessel function $I_1$, normalised so that $Z(0)=1$.

\begin{lemma}\label{q-lem:vmf}
For every $\kappa\ge0$ the law $\nu_b$ on the unit sphere $S^3$ satisfies
the Poincar\'e inequality $\Var_{\nu_b}(g)\le\frac13\int|\nabla g|^2d\nu_b$,
and the constant $\frac13$ is attained at $\kappa=0$.
\end{lemma}

\begin{proof}
In the polar angle $\theta\in(0,\pi)$ from $e_1$ the generator of the
reversible diffusion of $\nu_b$ is $L=\partial_\theta^2+b(\theta)
\partial_\theta+\sin^{-2}\theta\,\Delta_{S^2}$ with $b(\theta)=2\cot\theta-
\kappa\sin\theta$. On spherical harmonics of degree $\ell$ on $S^2$ it acts as
a Sturm--Liouville operator $L_\ell$ with potential
$\ell(\ell+1)\sin^{-2}\theta$. A ground state comparison of Barta's
type~\cite{Bar37} with the trial function $\psi=\sin\theta\,
e^{a\kappa\cos\theta}$ gives $-L_1\psi\ge3\psi$ for $a=-\frac15$, and the same
trial function with $a=-\frac25$ gives the bound $3$ for the derivatives of
zonal eigenfunctions; Appendix~\ref{app:vmfcomp} carries out both computations.
Hence every nonzero eigenvalue of $-L$ is at least $3$. At $\kappa=0$ the
first eigenvalue of the Laplacian of the unit $S^3$ is $3$.
\end{proof}

\subsection{Moments}

Let $A(\kappa)=\E_{\nu_b}x_1=Z'(\kappa)/Z(\kappa)=I_2(\kappa)/I_1(\kappa)$.

\begin{lemma}\label{q-lem:moments}
For every $\kappa>0$, in coordinates with $b=\kappa e_1$,
\begin{enumerate}[label=\textup{(\roman*)},leftmargin=*]
\item $\E x_2^2=\E x_3^2=\E x_4^2=A/\kappa$, $\E(1-x_1^2)=3A/\kappa$, the
covariance matrix of $x$ is $\operatorname{diag}(A',A/\kappa,A/\kappa,
A/\kappa)$, and $A'=1-A^2-3A/\kappa$;
\item $\displaystyle\frac{\kappa}{2+\sqrt{\kappa^2+4}}\le
A(\kappa)\le\frac{\kappa}{1+\sqrt{\kappa^2+9}}$;
\item $A'\le A/\kappa\le\bigl(1+\sqrt{\kappa^2+9}\bigr)^{-1}\le\frac14$;
\item $A(\kappa)\ge\kappa\big/\bigl(\frac32+\sqrt{\kappa^2+25/4}\bigr)$.
\end{enumerate}
\end{lemma}

\begin{proof}
(i) Integrating by parts against $\frac{d}{dt}(1-t^2)^{3/2}=-3t(1-t^2)^{1/2}$
gives $\kappa\,\E(1-x_1^2)=3\E x_1$, and the three coordinates
$x_2,x_3,x_4$ are exchangeable under $\nu_b$, with $\sum_ix_i^2=1$. The
off-diagonal covariances vanish by the reflections $x_i\mapsto-x_i$,
$i\ge2$, and $A'=\Var x_1$ because $A=(\log Z)'$.
(ii) By (i), $A$ solves the Riccati equation $A'=1-A^2-3A/\kappa$ with
$A(\kappa)=\kappa/4+O(\kappa^3)$. Appendix~\ref{app:vmfcomp} shows that the left
side of (ii) is a subsolution and the right side a supersolution of this
equation, with the same behaviour at $0$, and concludes by comparison.
(iii) The middle inequality is (ii). For the first, put $s=\sqrt{\kappa^2+4}$.
By (ii), $4A/\kappa\ge4/(2+s)$ and $1-A^2\le1-\kappa^2/(2+s)^2$, and
\[
\begin{aligned}
1-\frac{\kappa^2}{(2+s)^2}\le\frac4{2+s}
&\iff(2+s)^2-\kappa^2\le4(2+s)\\
&\iff s^2-\kappa^2\le4,
\end{aligned}
\]
which holds with equality. So $A'=1-A^2-3A/\kappa\le A/\kappa$.
(iv) is the comparison of (ii) with a third barrier, which improves the
lower bound of (ii); it is carried out in the same appendix.
\end{proof}

\subsection{The Lipschitz rate of the family}

Let $\delta$ be the chordal metric $\delta(x,x')=|x-x'|$ on $S^3$ and $W_1$
the Kantorovich distance for $\delta$. For $c\in\R^4$ and a unit vector
$h\in\R^4$ put
\begin{equation}\label{q-eq:rate}
L(c,h)=\sup_g\Cov_{\nu_c}\bigl(g,\ip hx\bigr),
\end{equation}
the supremum over $1$-Lipschitz functions $g$ for $\delta$. Since
$\frac{d}{dt}\nu_{c+th}(g)=\Cov_{\nu_{c+th}}(g,\ip hx)$, the
Kantorovich--Rubinstein duality gives, for all $b,b'$,
\begin{equation}\label{q-eq:w1segment}
W_1(\nu_b,\nu_{b'})\le|b-b'|\sup_{c\in[b,b']}L\Bigl(c,\frac{b'-b}{|b'-b|}
\Bigr).
\end{equation}
By rotation invariance $L(c,h)$ depends only on $\kappa=|c|$ and on the angle
$\varphi$ between $h$ and $c^\perp$, and we write $L(\kappa,\varphi)$.

\begin{lemma}\label{q-lem:rate}
For all $c$ and $h$, $L(c,h)\le1/(2\sqrt3)$. Moreover:
\begin{enumerate}[label=\textup{(\roman*)},leftmargin=*]
\item at $c=0$, $L(0,h)=\frac14$ for every $h$;
\item if $h\perp c$, then $L(c,h)=\E_{\nu_c}x_2^2=A(\kappa)/\kappa\le\frac14$;
\item if $\kappa\ge\sqrt{88}/3\approx3.127$, then $L(c,h)\le\frac14$ for
every $h$.
\end{enumerate}
\end{lemma}

\begin{proof}
A $1$-Lipschitz function for $\delta$ is $1$-Lipschitz for the geodesic
distance, which is larger, so $|\nabla g|\le1$ almost everywhere, and
Lemma~\ref{q-lem:vmf} gives $\Var g\le\frac13$. By Lemma~\ref{q-lem:moments}(i),
(iii) the largest eigenvalue of the covariance of $x$ is $A/\kappa\le\frac14$,
so $\Var\ip hx\le A/\kappa$. The Cauchy--Schwarz inequality gives
\begin{equation}\label{q-eq:csbound}
L(c,h)\le\bigl(\tfrac13A/\kappa\bigr)^{1/2}\le\Bigl(3\bigl(1+
\sqrt{\kappa^2+9}\bigr)\Bigr)^{-1/2}\le\frac1{2\sqrt3}.
\end{equation}
(i) and (ii) are reflection arguments. Let $R$ be the reflection in the
hyperplane $h^\perp$, and suppose $\nu_c$ is $R$-invariant, which is the case
when $c=0$ or $h\perp c$. Then $\E\ip hx=0$ and
\[
\begin{aligned}
\Cov(g,\ip hx)&=\E\bigl[(g(x)-g(Rx))\ip hx;\ \ip hx>0\bigr]\\
&\le\E\bigl[2\ip hx^2;\ \ip hx>0\bigr]=\E\ip hx^2,
\end{aligned}
\]
since $\delta(x,Rx)=2|\ip hx|$. Equality holds for $g=\ip hx$, which is
$1$-Lipschitz. At $c=0$, $\E\ip hx^2=\frac14$; for $h\perp c$, $\E\ip
hx^2=A/\kappa$. (iii) By~\eqref{q-eq:csbound}, $L\le\frac14$ as soon as
$1+\sqrt{\kappa^2+9}\ge\frac{16}3$, that is $\kappa^2\ge\frac{88}9$.
\end{proof}

The reflection settles the directions
$h\perp c$ and the point $c=0$; for an oblique $h$ the reflection no longer
preserves $\nu_c$, and the Cauchy--Schwarz bound~\eqref{q-eq:csbound} is too
weak for $\kappa<\sqrt{88}/3$. The next theorem closes this case.

\begin{theorem}[The sharp quarter]\label{q-thm:quarter}
For every $c\in\R^4$ and every unit vector $h\in\R^4$,
\[
L(c,h)\le\tfrac14 ,
\]
with equality at $c=0$. Equivalently, $W_1(\nu_b,\nu_{b'})\le\frac14|b-b'|$
for all $b,b'\in\R^4$, and the constant $\frac14$ is the best possible.
\end{theorem}

The equivalence is~\eqref{q-eq:w1segment} in one direction and the derivative
formula above it in the other; the sharpness is Lemma~\ref{q-lem:rate}(i). The
metric matters.\footnote{In the geodesic metric of $S^3$ the statement is
false already at $c=0$: the function $\arcsin x_2$ is geodesically
$1$-Lipschitz and its covariance with $x_2$ under the uniform law is
$\frac8{9\pi}\approx0.283>\frac14$. A proof therefore has to use that the
chordal distance between two points is shorter than their geodesic distance,
that is, it has to see the sphere from the ball it bounds. The proof below
does exactly this: it extends a chordally Lipschitz function to the ball and
integrates by parts there.} The proof has four steps. A Lipschitz function
on $S^3$ extends to the unit ball $B\subset\R^4$, and Green's formula in $B$
bounds the covariance by a negative Sobolev norm of the density
(Lemma~\ref{q-lem:harmonic}). This norm splits into a zonal and a transverse
part, both explicit series in modified Bessel functions
(Lemma~\ref{q-lem:spectral}). The transverse part is at most $\frac14$ because a
certain power series has nonnegative coefficients
(Proposition~\ref{q-prop:transverse}), and the zonal part is at most $\frac14$
by a Riccati comparison (Proposition~\ref{q-prop:zonal}).

\subsection{Harmonic extension}

Let $\sigma$ be the normalised surface measure of $S^3$ and $H_k$ the space
of spherical harmonics of degree $k$, the restrictions to $S^3$ of harmonic
homogeneous polynomials of degree $k$ on $\R^4$. The spaces $H_k$ are
orthogonal in $L^2(\sigma)$ and span it; we write $u_k$ for the component of
$u$ in $H_k$, and for $u$ with $\int u\,d\sigma=0$ we put
\begin{equation}\label{q-eq:Qnorm}
\mathcal Q(u)=\sum_{k\ge1}\frac1k\,\|u_k\|_{L^2(\sigma)}^2 .
\end{equation}
This is $\ip u{\Lambda^{-1}u}$ for the Dirichlet-to-Neumann operator
$\Lambda$ of $B$, which acts on $H_k$ as multiplication by $k$.

\begin{lemma}\label{q-lem:harmonic}
Let $u$ be a real-analytic function on $S^3$ with $\int u\,d\sigma=0$ and let
$g$ be $1$-Lipschitz for the chordal metric. Then
\[
\int_{S^3}g\,u\,d\sigma\le\tfrac12\,\mathcal Q(u)^{1/2}.
\]
\end{lemma}

\begin{proof}
The chordal metric is the Euclidean metric of $\R^4$ restricted to $S^3$,
so McShane's extension $G(y)=\inf_{x\in S^3}\bigl(g(x)+|x-y|\bigr)$
is $1$-Lipschitz on $\R^4$ and equals $g$ on $S^3$~\cite{McS34}; by
Rademacher's theorem $|\nabla G|\le1$ almost everywhere. Let
$\psi=\sum_{k\ge1}k^{-1}r^ku_k(\omega)$ in polar coordinates $y=r\omega$. Each
term is a harmonic polynomial, the series converges with all derivatives on
the closed ball because $u$ is real-analytic,\footnote{For real-analytic
$u$ the norms $\|u_k\|$ decay geometrically in $k$, and the supremum of a
function in $H_k$ is at most $(\dim H_k)^{1/2}=k+1$ times its $L^2(\sigma)$
norm, so the series for $\psi$ and for each of its derivatives converge
uniformly on $\bar B$.} and $\partial_r\psi=u$ on $S^3$. Green's formula for
the Lipschitz function $G$ and the smooth harmonic function $\psi$ reads
$\int_{S^3}G\,\partial_r\psi\,dS=\int_B\nabla G\cdot\nabla\psi\,dy$. With
$|S^3|=2\pi^2$ and $|B|=\pi^2/2$ the Cauchy--Schwarz inequality gives
\[
\int g\,u\,d\sigma=\frac1{2\pi^2}\int_B\nabla G\cdot\nabla\psi\,dy\le
\frac{|B|^{1/2}}{2\pi^2}\Bigl(\int_B|\nabla\psi|^2dy\Bigr)^{1/2}.
\]
Green's formula once more gives $\int_B|\nabla\psi|^2=\int_{S^3}\psi\,
\partial_r\psi\,dS=2\pi^2\sum_kk^{-1}\|u_k\|^2=2\pi^2\mathcal Q(u)$, and the
right side becomes $\frac12\mathcal Q(u)^{1/2}$.
\end{proof}

The lemma is a form of Thomson's principle. By Kantorovich--Rubinstein
duality, $\sup_g\int gu\,d\sigma$ is the least cost of a flow in $\R^4$ that
carries the positive part of $u\,d\sigma$ to its negative part, the cost
being the total length travelled by the mass.\footnote{This is the
Beckmann formulation of the transport problem with cost $|x-y|$: the
optimal cost equals the least $\int|V|$ over vector measures $V$ on $\R^4$
with divergence $u\,d\sigma$. Restricting the flow to the closed ball does
not change the infimum, since projecting onto $\bar B$ shortens paths. The
gradient $\nabla\psi$ on $B$ is one admissible flow. It minimises the energy
$\int|V|^2$ rather than the cost $\int|V|$, and the Cauchy--Schwarz inequality
converts one into the other at the price of the factor $|B|^{1/2}$.} The
harmonic field $\nabla\psi$ is one such flow, through the interior of the
ball, and the lemma computes its cost. At $c=0$ the flow is optimal: for
$u=\ip hx$ one has $\psi=\ip hy$, the field $\nabla\psi=h$ is constant,
the Cauchy--Schwarz inequality is an equality, and $\mathcal Q(u)=\|u\|^2=
\frac14$.

\subsection{Spectral decomposition}

Fix $c=\kappa e_1$ with $\kappa>0$ and write $h=\cos\varphi\,e_1+\sin\varphi
\,e_2$, which is no loss of generality by rotation invariance. Let
$\rho=d\nu_c/d\sigma=e^{\kappa x_1}/Z(\kappa)$, $Z(\kappa)=\int e^{\kappa x_1}
d\sigma=2I_1(\kappa)/\kappa$, and
\[
u=\bigl(\ip hx-\E_{\nu_c}\ip hx\bigr)\rho ,
\]
so that $\Cov_{\nu_c}(g,\ip hx)=\int gu\,d\sigma$ and $\int u\,d\sigma=0$.
Let $U_n$ be the Chebyshev polynomials of the second kind and put
\begin{equation}\label{q-eq:dn}
d_n(\kappa)=\frac{(n+1)\,I_{n+1}(\kappa)}{I_1(\kappa)},\qquad n\ge0 .
\end{equation}

\begin{lemma}\label{q-lem:spectral}
With this notation
\[
\mathcal Q(u)=\cos^2\varphi\,\mathcal Q_z(\kappa)+\sin^2\varphi\,
\mathcal Q_x(\kappa),
\]
where
\[
\mathcal Q_z=\sum_{n\ge1}\frac{d_n'(\kappa)^2}n,\qquad
\mathcal Q_x=\frac1{3\kappa^2}\sum_{n\ge1}(n+2)\,d_n(\kappa)^2 .
\]
\end{lemma}

\begin{proof}
The functions $U_n(x_1)$, $n\ge0$, are the zonal harmonics: $U_n(x_1)\in
H_n$, and they are orthonormal in $L^2(\sigma)$ (Appendix~\ref{app:bessel}).
The expansion $e^{\kappa t}=\frac2\kappa\sum_n(n+1)I_{n+1}(\kappa)U_n(t)$ of
the same appendix gives $\rho=\sum_nd_n(\kappa)U_n(x_1)$, with $d_0=1$.

Write $u=\cos\varphi\,u_z+\sin\varphi\,u_x$ with $u_z=(x_1-A)\rho$ and
$u_x=x_2\rho$; here $\E_{\nu_c}x_1=A$ and $\E_{\nu_c}x_2=0$. The reflection
$x_2\mapsto-x_2$ maps each $H_k$ to itself, fixes $u_z$ and reverses $u_x$,
so the components of $u_z$ and $u_x$ in each $H_k$ are orthogonal and the
cross term in $\mathcal Q(u)$ vanishes.

Since $A=Z'/Z$, we have $u_z=\partial_\kappa\rho=\sum_nd_n'(\kappa)U_n(x_1)$,
and the component of $u_z$ in $H_n$ is $d_n'U_n(x_1)$, of norm $|d_n'|$. This
gives $\mathcal Q_z$.

For $u_x$, let $y(t)=\cos t\,e_1+\sin t\,e_2$. The law $\nu_{\kappa y(t)}$ has
density $\rho_t=e^{\kappa\ip{y(t)}x}/Z(\kappa)=\sum_nd_n(\kappa)U_n(\ip{y(t)}x)$,
and $\partial_t\rho_t|_{t=0}=\kappa x_2\rho$. The function
$U_n(\ip{y(t)}x)$ lies in $H_n$ for every $t$, hence so does its derivative
$x_2U_n'(x_1)$ at $t=0$, and the component of $u_x$ in $H_n$ is
$\kappa^{-1}d_nx_2U_n'(x_1)$. Its norm follows from the eigenvalue
$n(n+2)$ of the Laplacian of $S^3$ on $H_n$. For $f=U_n(x_1)$,
\[
\begin{aligned}
n(n+2)&=n(n+2)\|f\|^2=\int|\nabla_{S^3}f|^2d\sigma\\
&=\int U_n'(x_1)^2(1-x_1^2)\,d\sigma ,
\end{aligned}
\]
since the tangential gradient of $x_1$ has squared length $1-x_1^2$. Given
$x_1$, the point $(x_2,x_3,x_4)$ is uniform on the sphere of radius
$(1-x_1^2)^{1/2}$, so $\E[x_2^2\mid x_1]=\frac13(1-x_1^2)$ and
$\|x_2U_n'(x_1)\|^2=\frac13n(n+2)$. Summing $n^{-1}\kappa^{-2}d_n^2\cdot\frac13
n(n+2)$ over $n\ge1$ gives $\mathcal Q_x$.
\end{proof}

By Lemmas~\ref{q-lem:harmonic} and~\ref{q-lem:spectral},
\begin{equation}\label{q-eq:Lbyq}
L(\kappa,\varphi)\le\tfrac12\max\bigl(\mathcal Q_z(\kappa),\mathcal Q_x(
\kappa)\bigr)^{1/2},
\end{equation}
and Theorem~\ref{q-thm:quarter} follows once both parts are at most
$\frac14$. Both are $\frac14$ at $\kappa=0$, by continuity from the values
$d_1\sim\kappa/2$, $d_1'\to\frac12$ and $d_n=O(\kappa^n)$.

\subsection{Three Bessel sums}

The series of Lemma~\ref{q-lem:spectral} reduce to closed forms. Write $I_j$
for $I_j(\kappa)$.

\begin{lemma}\label{q-lem:besselsums}
For every $\kappa>0$,
\begin{align*}
\sum_{j\ge1}j^2I_j^2&=\frac\kappa4\,I_1(2\kappa),\qquad
\sum_{j\ge1}j^3I_j^2=\frac{\kappa^2}4\bigl(I_0^2+I_1^2\bigr),\\
I_1(2\kappa)&=2\sum_{j\ge1}I_{j-1}I_j .
\end{align*}
\end{lemma}

The first and third identities follow from Parseval's formula for
$e^{\kappa\cos\theta}=\sum_{j\in\Z}I_je^{ij\theta}$, and the second from the
recurrences of $I_j$ by telescoping; Appendix~\ref{app:bessel} gives the
proofs.

\subsection{The transverse part}

\begin{proposition}\label{q-prop:transverse}
$\mathcal Q_x(\kappa)\le\frac14$ for every $\kappa\ge0$.
\end{proposition}

\begin{proof}
By~\eqref{q-eq:dn}, $\mathcal Q_x=(3\kappa^2I_1^2)^{-1}\sum_{j\ge2}j^2(j+1)
I_j^2$, and the sum equals $\sum_{j\ge1}(j^3+j^2)I_j^2-2I_1^2$. By
Lemma~\ref{q-lem:besselsums}, $\mathcal Q_x\le\frac14$ is equivalent to
\begin{equation}\label{q-eq:Dfunction}
\mathcal D(\kappa):=(8+2\kappa^2)I_1^2-\kappa^2I_0^2-\kappa I_1(2\kappa)\ge0 .
\end{equation}
Each term is an entire function of $\kappa$ that is even after the
factor $\kappa^2$ is removed, so $\mathcal D(\kappa)=\sum_{n\ge0}\mathcal D_n
\kappa^{2n}$. Lemma~\ref{q-lem:Dcoef} of Appendix~\ref{app:bessel} computes the
coefficients: $\mathcal D_0=\mathcal D_1=0$, and for $n=m+1\ge2$
\[
\mathcal D_n=\frac{N_m}{4^m\,m!\,(m+1)!\,(m+1)(m+2)},
\]
with
\[
\begin{aligned}
N_m={}&\binom{2m}m\bigl(4(2m+1)+(m-1)(m+1)(m+2)\bigr)\\
&-4^m(m+1)(m+2).
\end{aligned}
\]
The numerator $N_m$ vanishes for $m=0,1,2$ and is positive for $m\ge3$
(Lemma~\ref{q-lem:Dcoef}). So every coefficient of $\mathcal D$ is nonnegative
and $\mathcal D\ge0$.
\end{proof}

The first nonzero coefficient is $\mathcal D_4=\frac1{2304}$, so
$\frac14-\mathcal Q_x=\mathcal D/(12\kappa^2I_1^2)$ is of order $\kappa^4$ at $0$: the transverse part
leaves the value $\frac14$ very slowly, which is why no estimate that loses a
constant can prove the proposition.

\subsection{The zonal part}

Let $A=I_2/I_1$ as before, so that $d_1=2A$ and, by
Lemma~\ref{q-lem:moments}(i), $A'=1-A^2-3A/\kappa$.

\begin{proposition}\label{q-prop:zonal}
For $0<\kappa\le\sqrt{88}/3$,
\[
\mathcal Q_z(\kappa)\le\frac14+4A'^2+\frac{A^2}4-\frac A{2\kappa}-
\frac{2A^2}{\kappa^2}<\frac14 .
\]
\end{proposition}

\begin{proof}
\emph{The higher modes.} With $r_\nu=I_{\nu+1}/I_\nu$ the recurrence
$I_\nu'=I_{\nu+1}+\frac\nu\kappa I_\nu$ gives
\[
\frac{d_n'}{d_n}=\frac{I_{n+1}'}{I_{n+1}}-\frac{I_1'}{I_1}=\frac n\kappa+
r_{n+1}-r_1 .
\]
By Lemma~\ref{q-lem:ratio} of Appendix~\ref{app:bessel}, $r_{n+1}\le r_1$, and
$r_1=A\le\kappa/4$ by Lemma~\ref{q-lem:moments}(iii). Hence $-n/\kappa\le
n/\kappa-\kappa/4\le d_n'/d_n\le n/\kappa$ for $n\ge2$ and $\kappa^2\le16$, so
$d_n'^2\le n^2d_n^2/\kappa^2$. With $d_1'=2A'$,
\[
\mathcal Q_z\le4A'^2+\frac1{\kappa^2}\sum_{n\ge2}n\,d_n^2=4A'^2+\frac{S_3-
S_2-4I_2^2}{\kappa^2I_1^2},
\]
where $S_2=\sum_{j\ge1}j^2I_j^2$ and $S_3=\sum_{j\ge1}j^3I_j^2$; the
second form uses $n\,d_n^2I_1^2=(j-1)j^2I_j^2$ with $j=n+1$.

\emph{The closed form.} The recurrences $I_0-I_2=\frac2\kappa I_1$ and
$I_1-I_3=\frac4\kappa I_2$ give $I_0/I_1=2/\kappa+A$ and $I_2I_3/I_1^2=A-
4A^2/\kappa$. By Lemma~\ref{q-lem:besselsums},
$S_3/(\kappa^2I_1^2)=\frac14(2/\kappa+A)^2+\frac14$, and, keeping three
terms of the third identity,
\[
\frac{S_2}{\kappa^2I_1^2}\ge\frac{I_0I_1+I_1I_2+I_2I_3}{2\kappa I_1^2}=
\frac1{2\kappa}\Bigl(\frac2\kappa+3A-\frac{4A^2}\kappa\Bigr).
\]
Collecting terms, $\mathcal Q_z\le\frac14+4A'^2+T(A)$ with $T(a)=\frac{a^2}4-
\frac a{2\kappa}-\frac{2a^2}{\kappa^2}$, which is the first inequality.

\emph{The barrier.} Let $s=\sqrt{\kappa^2+25/4}$ and $A_\flat=\kappa/(\frac32
+s)$. By Lemma~\ref{q-lem:moments}(iv), $A_\flat\le A\le\kappa/4$. The function
$a\mapsto1-a^2-3a/\kappa$ decreases on $a\ge0$, and $A'\ge0$ because $A'=
\Var x_1$, so $0\le A'\le\alpha:=1-A_\flat^2-3A_\flat/\kappa$. For $\kappa^2\le
12$ the quadratic $T$ decreases on $[0,\kappa/4]$, since
\[
T(a)-T(a')=(a'-a)\,\frac{2\kappa-(a+a')(\kappa^2-8)}{4\kappa^2}
\]
and $(a+a')(\kappa^2-8)\le\frac\kappa2(\kappa^2-8)\le2\kappa$ there. So
$\mathcal Q_z\le\frac14+4\alpha^2+T(A_\flat)$. With $\kappa^2=s^2-\frac{25}4$
and $A_\flat/\kappa=2/(2s+3)$ both terms are rational in $s$:
$\alpha=16/(2s+3)^2$, and
\begin{equation}\label{q-eq:zonalfactor}
\begin{aligned}
4\alpha^2+T(A_\flat)&=\frac{(2s-5)\,p(s)}{4(2s+3)^4},\\
p(s)&=8s^3+28s^2-98s-695 .
\end{aligned}
\end{equation}
For $0<\kappa\le\sqrt{88}/3$ we have $\frac52<s\le\sqrt{577}/6<4.004$. On
$s\ge\frac52$ the cubic $p$ increases, since $p'(s)=24s^2+56s-98>0$, and
$p(4.004)<-120$. So the right side of~\eqref{q-eq:zonalfactor} is negative.
\end{proof}

\subsection{Proof of Theorem~\ref{q-thm:quarter}}

At $c=0$ the theorem is Lemma~\ref{q-lem:rate}(i), and for $|c|\ge\sqrt{88}/3$
it is Lemma~\ref{q-lem:rate}(iii). For $0<|c|<\sqrt{88}/3$,
Propositions~\ref{q-prop:transverse} and~\ref{q-prop:zonal} give $\mathcal Q_x,
\mathcal Q_z\le\frac14$, and~\eqref{q-eq:Lbyq} gives $L\le\frac14$. \qed


\section{Computations for the single-link law}\label{app:vmfcomp}

\subsection{Barta's inequality}

Let $\rho(\theta)=\sin^2\theta\,e^{\kappa\cos\theta}$ on $(0,\pi)$, let $V$
be a continuous potential, and let
\[
T=-\rho^{-1}\partial_\theta\,\rho\,\partial_\theta+V .
\] If $\psi>0$ is smooth on $(0,\pi)$ with
$\rho\psi^2=O(\sin^4\theta)$ at the end points and $T\psi\ge\lambda_0\psi$,
then
\[
\int(|h'|^2+Vh^2)\rho\,d\theta\ge\lambda_0\int h^2\rho\,d\theta
\]
for every smooth $h$ vanishing to first order at $0$ and $\pi$. Indeed, writing
$h=\psi u$ and integrating by parts, the boundary terms $\rho\psi\psi'u^2$
vanish and
\[
\int(|h'|^2+Vh^2)\rho=\int\rho\psi^2|u'|^2+\int\frac{T\psi}\psi\,h^2\rho
\ \ge\ \lambda_0\int h^2\rho .
\]
This is the inequality of Barta~\cite{Bar37} in the form of a ground state
transformation.

\subsection{The computations of Lemma~\ref{q-lem:vmf}}

The generator of the reversible diffusion of $\nu_b$ on $S^3$ is
$L=\Delta_{S^3}+\kappa\nabla x_1\cdot\nabla$. In the polar angle $\theta$
from $e_1$ the round metric is $d\theta^2+\sin^2\theta\,g_{S^2}$, so
$\Delta_{S^3}=\partial_\theta^2+2\cot\theta\,\partial_\theta+\sin^{-2}\theta\,
\Delta_{S^2}$, and $\nabla x_1\cdot\nabla=-\sin\theta\,\partial_\theta$. On
$h(\theta)Y_\ell$ with a spherical harmonic $Y_\ell$ of degree $\ell$ on $S^2$
the operator acts as $L_\ell=\partial_\theta^2+b\,\partial_\theta-\ell(\ell+1)
\sin^{-2}\theta$, $b=2\cot\theta-\kappa\sin\theta=\rho'/\rho$, so
$-L_\ell=T$ with $V=\ell(\ell+1)\sin^{-2}\theta$.

\emph{Degree $\ell\ge1$.} Take $\psi=\sin\theta\,e^{a\kappa\cos\theta}$. Then
$\psi'/\psi=\cot\theta-a\kappa\sin\theta$ and
$\psi''/\psi=(\psi'/\psi)^2+(\psi'/\psi)'$, and a short computation gives
\[
\frac{-L_1\psi}\psi=3+(1+5a)\kappa\cos\theta-a(1+a)\kappa^2\sin^2\theta .
\]
For $a=-\frac15$ this is $3+\frac4{25}\kappa^2\sin^2\theta\ge3$. Since
$V=\ell(\ell+1)\sin^{-2}\theta\ge2\sin^{-2}\theta$ for $\ell\ge1$, Barta's
inequality gives $-L_\ell\ge3$ on every sector with $\ell\ge1$.

\emph{Degree $\ell=0$.} Let $g(\theta)$ be a zonal eigenfunction,
$-L_0g=\lambda g$, $g$ not constant, so that $g'(0)=g'(\pi)=0$. Differentiating
$g''+bg'=-\lambda g$ gives, for $h=g'$,
\[
-\bigl(h''+bh'\bigr)-b'h=\lambda h,\qquad -b'=\frac2{\sin^2\theta}+\kappa
\cos\theta,
\]
so $h$ is an eigenfunction of $T$ with $V=2\sin^{-2}\theta+\kappa\cos\theta$,
with eigenvalue $\lambda$, and $h\not\equiv0$ vanishes at both ends. With the
same $\psi$,
\[
\frac{T\psi}\psi=3+(2+5a)\kappa\cos\theta-a(1+a)\kappa^2\sin^2\theta,
\]
which for $a=-\frac25$ is $3+\frac6{25}\kappa^2\sin^2\theta\ge3$. Barta's
inequality applied to $h$ gives $\lambda\ge3$. Together with the case
$\ell\ge1$ every nonzero eigenvalue of $-L$ is at least $3$, which is
Lemma~\ref{q-lem:vmf}; the spectrum is discrete because $S^3$ is compact and
the drift is smooth.

\subsection{Moments}

The density of $t=x_1$ under $\nu_b$ is proportional to $(1-t^2)^{1/2}e^{
\kappa t}$. For every smooth $\phi$ on $[-1,1]$,
\[
\frac d{dt}\Bigl[(1-t^2)^{3/2}e^{\kappa t}\Bigr]=(1-t^2)^{1/2}e^{\kappa t}
\bigl(\kappa(1-t^2)-3t\bigr),
\]
and integration by parts gives $\E[\phi'(t)(1-t^2)]=\E\bigl[\phi(t)\bigl(3t-
\kappa(1-t^2)\bigr)\bigr]$. With $\phi=1$ this is $\kappa\E(1-x_1^2)=3\E x_1$,
the identity of Lemma~\ref{q-lem:moments}(i). With $\phi=t^k$ it gives the
recursion
\[
m_{k+2}=m_k-\frac{k+3}\kappa\,m_{k+1}+\frac k\kappa\,m_{k-1},\qquad m_k=\E
x_1^k,
\]
from $m_0=1$ and $m_1=A$; all polynomial moments of $\nu_b$ are rational in
$\kappa$ and $A(\kappa)$.\footnote{For instance $m_2=1-3A/\kappa$,
$m_3=A-3/\kappa+12A/\kappa^2$ and $m_4=m_2-5m_3/\kappa+2A/\kappa$. The joint moments
with $x_2,x_3,x_4$ follow from exchangeability:
$\E[\phi(x_1)x_2^2]=\frac13\E[\phi(x_1)(1-x_1^2)]$.}

\subsection{The Riccati barriers}

By Lemma~\ref{q-lem:moments}(i), $A$ solves $A'=\mathcal F(\kappa,A):=1-A^2-
3A/\kappa$ on $(0,\infty)$, and $A(\kappa)=\frac\kappa4-\frac{\kappa^3}{96}+
O(\kappa^5)$ from the series of $I_1$ and $I_2$. Let $s=\sqrt{\kappa^2+4}$,
$u=\sqrt{\kappa^2+9}$, $\phi_-=\kappa/(2+s)$ and $\phi_+=\kappa/(1+u)$. Direct
differentiation gives
\[
\begin{aligned}
\phi_-'-\mathcal F(\kappa,\phi_-)&=-\frac{\kappa^2}{s(s+2)^2}\le0,\\
\phi_+'-\mathcal F(\kappa,\phi_+)&=\frac{(u-3)^2}{u(1+u)^2}\ge0,
\end{aligned}
\]
and $\phi_\pm=\frac\kappa4+O(\kappa^3)$.\footnote{For the first identity write
$\phi_-=(s-2)/\kappa$, which follows from $s^2-4=\kappa^2$; then
$\phi_-'=(\kappa^2/s-s+2)/\kappa^2$ and $1-\phi_-^2-3\phi_-/\kappa=
(\kappa^2-(s-2)^2-3(s-2))/\kappa^2$, and the difference simplifies with
$s^2=\kappa^2+4$. The second is the same computation with
$\phi_+=\kappa(u-1)/(u^2-1)=\kappa(u-1)/(\kappa^2+8)$ in place of $\phi_-$
and $u^2=\kappa^2+9$ in place of $s^2=\kappa^2+4$. After clearing the denominator $\kappa^2(1+u)^2$ the
numerator becomes $\kappa^2(u-3)^2/u$, which is the stated residual.} The comparison is singular at $0$
because of the term $3A/\kappa$, and we carry it out with an integrating
factor. Let $D=A-\phi_-$. Then
\[
\begin{aligned}
D'&=\mathcal F(\kappa,A)-\mathcal F(\kappa,\phi_-)-\bigl(\phi_-'-\mathcal F(
\kappa,\phi_-)\bigr)\\
&=-\Bigl(A+\phi_-+\frac3\kappa\Bigr)D+R,
\end{aligned}
\]
with $R=\kappa^2/(s(s+2)^2)\ge0$. Put $\mu(\kappa)=\kappa^3\exp\int_0^\kappa
(A+\phi_-)$. Then $(\mu D)'=\mu R\ge0$, and $\mu D=O(\kappa^6)\to0$ as
$\kappa\to0$, so $\mu D\ge0$ and $A\ge\phi_-$. The same argument with
$D=\phi_+-A$ and $R=(u-3)^2/(u(1+u)^2)$ gives $A\le\phi_+$. This is
Lemma~\ref{q-lem:moments}(ii).

For (iv) let $s=\sqrt{\kappa^2+25/4}$ and $\phi_\flat=\kappa/(\frac32+s)$.
Since $s^2-\frac94=\kappa^2+4$, the same differentiation gives
\[
\phi_\flat'-\mathcal F(\kappa,\phi_\flat)=-\frac{5\,(s-\frac52)}{2\,s\,
(\frac32+s)^2}\le0,
\]
so $\phi_\flat$ is a subsolution, and $\phi_\flat=\frac\kappa4-\frac{\kappa^3}
{80}+O(\kappa^5)$.\footnote{With $\phi_\flat=\kappa(s-\frac32)/(\kappa^2+4)$ one
has $\phi_\flat'=(\frac32 s+\frac{25}4)/(s(\frac32+s)^2)$, using $ds/d\kappa=
\kappa/s$; and $\mathcal F(\kappa,\phi_\flat)=1-\phi_\flat^2-3/(\frac32+s)$.
Clearing the denominator $s(\frac32+s)^2$ and replacing $\kappa^2$ by
$s^2-\frac{25}4$ leaves $-\frac52 s+\frac{25}4$ in the numerator, which is the
stated residual. The expansion follows from $s=\frac52+\frac{\kappa^2}5+
O(\kappa^4)$.} Hence $A-\phi_\flat=\frac{\kappa^3}{480}+O(\kappa^5)$, and the
integrating-factor argument above, with $D=A-\phi_\flat$, $R=\frac{5(s-5/2)}
{2s(3/2+s)^2}\ge0$ and $\mu=\kappa^3\exp\int_0^\kappa(A+\phi_\flat)$, gives
$\mu D=O(\kappa^6)$ at $0$, $(\mu D)'=\mu R\ge0$, and so $A\ge\phi_\flat$. Since
$\frac32+\sqrt{\kappa^2+25/4}<2+\sqrt{\kappa^2+4}$ for $\kappa>0$, this
improves the lower bound of (ii).


\section{Bessel functions}\label{app:bessel}

This appendix proves the auxiliary statements used in Appendix~\ref{app:sharp}
for Theorem~\ref{q-thm:quarter}. Throughout, $I_j=I_j(\kappa)$ with $\kappa>0$,
and $I_{-j}=I_j$ for integers $j$. We use the series
\[
I_\nu(\kappa)=\sum_{i\ge0}\frac{(\kappa/2)^{2i+\nu}}{i!\,(i+\nu)!},
\]
the integral $I_j(\kappa)=\pi^{-1}\int_0^\pi e^{\kappa\cos\theta}\cos
j\theta\,d\theta$, the Fourier series $e^{\kappa\cos\theta}=\sum_{j\in\Z}
I_je^{ij\theta}$ and the recurrences
\begin{equation}\label{q-eq:besselrec}
I_{j-1}-I_{j+1}=\frac{2j}\kappa\,I_j,\quad I_j'=I_{j+1}+\frac j\kappa\,I_j ,
\end{equation}
all of which are in Watson's treatise~\cite{Wat44}, \S\S3.7, 3.71,
2.22. All series below converge absolutely, because $0<I_j\le(\kappa/2)^j
e^\kappa/j!$.

\subsection{Zonal harmonics}

Under $\sigma$ the coordinate $x_1$ has the density $\frac2\pi(1-t^2)^{1/2}$
on $[-1,1]$, the area of the two-sphere $\{x_1=t\}$ being proportional to
$1-t^2$ and its normal speed to $(1-t^2)^{-1/2}$. The Chebyshev polynomials
$U_n(\cos\theta)=\sin((n+1)\theta)/\sin\theta$ satisfy
\[
\int U_nU_m\,d\sigma=\frac2\pi\int_0^\pi\sin((n+1)\theta)\sin((m+1)\theta)
\,d\theta=\delta_{nm},
\]
and $U_n=C_n^{(1)}$ is the Gegenbauer polynomial of index $\frac{4-2}2=1$, so
$U_n(\ip yx)$ is the zonal harmonic of degree $n$ on $S^3$ with pole $y$:
$|x|^nU_n(\ip yx/|x|)$ is a harmonic polynomial on $\R^4$~\cite{SW71},
Ch.~IV, \S2.

\begin{lemma}\label{q-lem:chebexp}
$e^{\kappa t}=\frac2\kappa\sum_{n\ge0}(n+1)I_{n+1}(\kappa)\,U_n(t)$ in
$L^2\bigl(\frac2\pi(1-t^2)^{1/2}dt\bigr)$, and uniformly on $[-1,1]$.
\end{lemma}

\begin{proof}
With $t=\cos\theta$ and $2\sin\theta\sin((n+1)\theta)=\cos n\theta-\cos(
n+2)\theta$, the $n$-th coefficient is
\[
\begin{aligned}
\frac2\pi\int_0^\pi e^{\kappa\cos\theta}\sin((n+1)\theta)\sin\theta\,d\theta
&=I_n-I_{n+2}\\
&=\frac{2(n+1)}\kappa I_{n+1},
\end{aligned}
\]
by~\eqref{q-eq:besselrec}. Since $|U_n|\le n+1$ and $I_{n+1}$ decays
factorially, the series converges uniformly.
\end{proof}

\subsection{Proof of Lemma~\ref{q-lem:besselsums}}

Let $F(\theta)=e^{\kappa\cos\theta}=\sum_jI_je^{ij\theta}$. Then
\[
F'(\theta)=-\kappa\sin\theta\,F(\theta)=\sum_jijI_je^{ij\theta},
\]
and Parseval's formula gives
\[
\begin{aligned}
\sum_{j\in\Z}j^2I_j^2&=\frac{\kappa^2}{2\pi}\int_0^{2\pi}\sin^2\theta\,
e^{2\kappa\cos\theta}d\theta\\
&=\frac{\kappa^2}2\bigl(I_0(2\kappa)-I_2(2\kappa)\bigr)=\frac\kappa2I_1(2\kappa),
\end{aligned}
\]
the last step by~\eqref{q-eq:besselrec} at the argument $2\kappa$. Halving gives
the first identity. The coefficient of $e^{i\theta}$ in $F^2=e^{2\kappa\cos
\theta}$ is $I_1(2\kappa)=\sum_{j\in\Z}I_jI_{1-j}$, and the terms with $j\le0$
repeat those with $j\ge1$; this is the third identity.

For the second, $jI_j=\frac\kappa2(I_{j-1}-I_{j+1})$ by~\eqref{q-eq:besselrec}, so
\[
\sum_{j\ge1}j^3I_j^2=\frac{\kappa^2}4\sum_{j\ge1}j\bigl(I_{j-1}^2+I_{j+1}^2-
2I_{j-1}I_{j+1}\bigr).
\]
Shifting indices gives
\[
\sum_{j\ge1}j\bigl(I_{j-1}^2+I_{j+1}^2\bigr)=I_0^2+\sum_{m\ge1}2mI_m^2,
\qquad
\sum_{m\ge1}2mI_m^2=\kappa\sum_{m\ge1}(I_{m-1}I_m-I_mI_{m+1})=\kappa I_0I_1,
\]
the last by telescoping. For the mixed term write $j=\frac12((j-1)+(j+1))$
and apply~\eqref{q-eq:besselrec} to $(j-1)I_{j-1}$ and to $(j+1)I_{j+1}$,
using $I_{-1}=I_1$ at $j=1$:
\[
\begin{aligned}
jI_{j-1}I_{j+1}=\frac\kappa4\bigl[&(I_{j-1}I_j-I_jI_{j+1})\\
&+(I_{j-2}I_{j+1}-I_{j-1}I_{j+2})\bigr].
\end{aligned}
\]
Both brackets telescope, to $I_0I_1$ and to $I_{-1}I_2=I_1I_2$, so
$\sum_{j\ge1}jI_{j-1}I_{j+1}=\frac\kappa4(I_0I_1+I_1I_2)$. Altogether
\[
\sum_{j\ge1}j^3I_j^2=\frac{\kappa^2}4\Bigl(I_0^2+\frac\kappa2I_1(I_0-I_2)
\Bigr)=\frac{\kappa^2}4\bigl(I_0^2+I_1^2\bigr). \qquad\qed
\]

\subsection{\texorpdfstring{The coefficients of $\mathcal D$}{The coefficients of D}}

\begin{lemma}\label{q-lem:Dcoef}
Write $\mathcal D(\kappa)=\sum_{n\ge0}\mathcal D_n\kappa^{2n}$ for the function
of~\eqref{q-eq:Dfunction}. Then $\mathcal D_0=\mathcal D_1=0$, and for $m\ge1$
\[
\mathcal D_{m+1}=\frac{N_m}{4^m\,m!\,(m+1)!\,(m+1)(m+2)},
\]
with $N_m=\binom{2m}m\bigl(4(2m+1)+(m-1)(m+1)(m+2)\bigr)-4^m(m+1)(m+2)$.
Moreover $N_m=0$ for $m\le2$ and $N_m>0$ for $m\ge3$.
\end{lemma}

\begin{proof}
\emph{The coefficients.} Multiplying the series of $I_0$ and $I_1$ and
collecting powers,
\[
\begin{gathered}
I_0^2=\sum_{m\ge0}\alpha_m\kappa^{2m},\qquad I_1^2=\sum_{m\ge0}\beta_m
\kappa^{2m+2},\\
\kappa I_1(2\kappa)=\sum_{m\ge0}\frac{\kappa^{2m+2}}{m!\,(m+1)!},
\end{gathered}
\]
where
\[
\begin{aligned}
\alpha_m&=\frac1{4^m}\sum_{i=0}^m\frac1{i!^2(m-i)!^2},\\
\beta_m&=\frac1{4^{m+1}}\sum_{i=0}^m\frac1{i!\,(i+1)!\,(m-i)!\,(m-i+1)!}.
\end{aligned}
\]
By Vandermonde's identity, $\sum_i\binom mi^2=\binom{2m}m$ and
$\sum_i\binom mi\binom{m+2}{i+1}=\binom{2m+2}{m+1}$, so
\[
\alpha_m=\frac{\binom{2m}m}{4^m\,m!^2},\qquad\beta_m=\frac{\binom{2m+2}{m+1}}
{4^{m+1}m!\,(m+2)!}.
\]
The coefficient of $\kappa^{2m+2}$ in $\mathcal D$ is $\mathcal D_{m+1}=8\beta_m
+2\beta_{m-1}-\alpha_m-1/(m!(m+1)!)$, with $\beta_{-1}=0$; for $m=0$ it is
$\frac84-1-1=0$, and $\mathcal D_0=0$ because every term of $\mathcal D$ is
$O(\kappa^2)$. For $m\ge1$ put $\Delta=4^m\,m!\,(m+1)!\,(m+1)(m+2)$ and
$b=\binom{2m}m$. With $\binom{2m+2}{m+1}=\frac{2(2m+1)}{m+1}b$,
\[
\begin{aligned}
8\beta_m&=\frac{4(2m+1)\,b}\Delta,\qquad\alpha_m=\frac{(m+1)^2(m+2)\,b}\Delta,\\
2\beta_{m-1}&=\frac{2m(m+1)(m+2)\,b}\Delta,\\
\frac1{m!(m+1)!}&=\frac{4^m(m+1)(m+2)}\Delta,
\end{aligned}
\]
and $2m-(m+1)=m-1$ gives $\mathcal D_{m+1}=N_m/\Delta$.

\emph{The sign.} Directly, $N_0=N_1=N_2=0$ ($2-2$, $24-24$, $192-192$),
$N_3=80$, $N_4=1140$ and $N_5=10416$. For $m\ge6$ we first show
$16^m\le4m\binom{2m}m^2$ for $m\ge1$ by induction: it is an equality at
$m=1$, and $(m+1)\binom{2m+2}{m+1}=2(2m+1)\binom{2m}m$ gives
\[
\begin{aligned}
4(m+1)\binom{2m+2}{m+1}^2&=\frac{16(2m+1)^2}{m+1}\binom{2m}m^2\\
&\ge\frac{16(2m+1)^2}{4m(m+1)}16^m\ge16^{m+1}.
\end{aligned}
\]
Since $(m-1)^2\ge4m$ for $m\ge6$, it follows that $\binom{2m}m^2(m-1)^2\ge
16^m$, that is $\binom{2m}m(m-1)\ge4^m$, and
\[
\begin{aligned}
N_m={}&\Bigl(\binom{2m}m(m-1)-4^m\Bigr)(m+1)(m+2)\\
&+4(2m+1)\binom{2m}m>0 .
\end{aligned}
\]
\end{proof}

\subsection{Monotonicity of Bessel ratios}

\begin{lemma}\label{q-lem:ratio}
For $n\ge1$ and $\kappa>0$, $I_{n+2}(\kappa)/I_{n+1}(\kappa)\le I_2(\kappa)/
I_1(\kappa)$.
\end{lemma}

\begin{proof}
We use a lemma of Biernacki and Krzy\.z~\cite{BK55}: if $f=\sum a_jx^j$ and
$g=\sum b_jx^j$ have positive coefficients and $a_j/b_j$ is nonincreasing in
$j$, then $f/g$ is nonincreasing on $x>0$. Indeed
\[
\begin{aligned}
f'g-fg'&=\sum_{i,j}(i-j)a_ib_jx^{i+j-1}\\
&=\frac12\sum_{i,j}(i-j)b_ib_j\Bigl(\frac{a_i}{b_i}-\frac{a_j}{b_j}\Bigr)
x^{i+j-1}\le0 .
\end{aligned}
\]
Take $x=\kappa^2$, $f=I_{n+1}(\kappa)/\kappa^{n+1}$ and $g=I_1(\kappa)/\kappa$.
Then $a_j/b_j=2^{-n}(j+1)!/(j+n+1)!$ decreases in $j$, so $I_{n+1}/(\kappa^n
I_1)$ is nonincreasing in $\kappa$. Its logarithmic derivative is, by the
second recurrence of~\eqref{q-eq:besselrec},
\[
\begin{aligned}
&\frac{n+1}\kappa+\frac{I_{n+2}}{I_{n+1}}-\frac n\kappa-\frac1\kappa-\frac{I_2}
{I_1}\\
&\qquad=\frac{I_{n+2}}{I_{n+1}}-\frac{I_2}{I_1}\le0 .
\end{aligned}\qedhere
\]
\end{proof}

